Durante o desenvolvimento deste trabalho, ferramentas de inteligência artificial generativa foram utilizadas como apoio em etapas de implementação, estruturação do sistema de testes, revisão de código e revisão textual.

Essas ferramentas também foram testadas em discussões de arquitetura, mas suas respostas não foram adotadas como base técnica do projeto. No contexto deste trabalho, elas se mostraram limitadas para decisões de sistemas embarcados, especialmente por sugerirem alternativas com novos pontos de falha ou por não considerarem recursos disponíveis na plataforma. Por isso, a arquitetura final do WCAN, as decisões de implementação, a interpretação dos resultados e a redação final foram revisadas e assumidas pelo autor. A inteligência artificial foi usada como ferramenta de apoio, não como fonte de validação técnica. O autor se responsabiliza por todo o conteúdo do documento.
